\section{Detailed Boundary Cases and Reviewer Checklist}
\label{app:boundaries}

\subsection{Boundary Cases}

\paragraph{Time conditioning versus action conditioning.} MRI CEKWorld predicts contrast-enhancement appearance at a selected elapsed time \citep{Kong2026MRIContrastEnhancement}. The model implements a directed temporal transition and is therefore Level~2. Time is not an executable intervention, so the system does not reach Level~3 unless a contrast dose, agent, or acquisition action is exposed and validated.

\paragraph{Care-process prediction versus patient simulation.} Foresight and related EHR generators predict longitudinal clinical events \citep{Kraljevi_2024Foresightgenerativepretrained,Waxler2025GenerativeMedicalEvent}. A medication or procedure emitted as the next event remains part of the predicted record. The model reaches Level~3 only when an external action can be set for the same history and the subsequent patient state is evaluated.

\paragraph{Treatment conditioning versus treatment effect.} Treatment-aware tumor and longitudinal imaging models can reach Level~3 when treatment changes the generated future \citep{Yang2025MedicalWorldModel,Liu2025TreatmentawareDiffusion}. They reach Level~4 only when the comparison between alternatives is supported by an explicit estimand and identification, structural, trial, or sensitivity argument.

\paragraph{Planning versus clinical actionability.} CLARITY, patient-dynamics agents, and model-based sepsis RL operationally use simulated trajectories for decisions and therefore exemplify Level~5 capability \citep{Ding2025CLARITYMedicalWorld,Wu2026AgentifyingPatientDynamics,Raghu2018ModelBasedReinforcement}. The label does not establish that the simulated action response is causal, that the reward is patient-centered, or that the policy is ready for clinical use.

\paragraph{Embodied action versus outcome validity.} Surgical tool motion, probe movement, and acquisition rotation provide clear executable actions \citep{Sivakumar2026SWoMoNeuroSymbolic,Yue2025EchoWorldLearningMotion,Yang2025Xray2XrayWorldModel}. They satisfy the Level~3 task-form boundary when they alter the modeled future, but visual or control success does not substitute for physical safety or clinical outcome validation.

\subsection{Checklist for Authors and Reviewers}

\begin{enumerate}[leftmargin=*]
\item \textbf{State:} What patient or environment state is represented, at what time scale, and which clinically important variables may be missing?
\item \textbf{Action:} Can the input be deliberately set for the same pre-action state? Is it an executable intervention, acquisition, or control, or only time, context, or an observed care event?
\item \textbf{Transition:} Where does the action enter, what variables can it affect, and is the model learned, mechanistic, or hybrid?
\item \textbf{Outcome:} Is the target a future observation, surrogate, treatment response, control success, or patient-relevant endpoint? Does it match the stated use?
\item \textbf{Validation:} Are horizon, support, uncertainty, external shift, causal assumptions, and decision consequences evaluated at the level required by the claim?
\item \textbf{Capability language:} Does the manuscript distinguish representation, future prediction, action-conditioned simulation, counterfactual reasoning, and planning rather than using them interchangeably?
\end{enumerate}

\begin{table*}[t]
\caption{Representative boundary cases spanning the existence gate and Level~1--5. Outside-scope examples are retained because branding or optimization language alone does not establish a patient or environment transition.}
\label{tab:core-matrix}
\centering
\scriptsize
\setlength{\tabcolsep}{3.0pt}
\renewcommand{\arraystretch}{1.05}
\begin{tabularx}{\textwidth}{@{}p{3.1cm}cp{2.4cm}p{3.6cm}Xp{1.1cm}@{}}
\toprule
\textbf{Paper} & \textbf{Level} & \textbf{Setting} & \textbf{Action contract} & \textbf{Transition contract} & \textbf{Code} \\
\midrule
CheXWorld: Exploring Image World Modeling for Radiograph Representation ... & 1 & medical imaging & Mask locations and image/domain transformation parameters are training conditions, ... & JEPA-style latent prediction across spatial context and image transformations; no future ... & \cgood yes \\
Enhancing Amyloid PET Quantification: MRI-Guided Super-Resolution Using Latent ... & 1 & medical imaging & No settable clinical action enters the model. & An MRI-guided latent diffusion model reconstructs a higher-resolution PET image for ... & \cgood yes \\
MRI Contrast Enhancement Kinetics World Model & 2 & medical imaging & Elapsed time is a query coordinate; contrast dose, agent, ... & A conditional generative model predicts contrast-enhanced MRI at arbitrary future time ... & \cgood yes \\
Forecasting individualized progression of Alzheimer's disease using structural ... & 2 & longitudinal progression & No treatment, acquisition, or clinical action is supplied; target ... & A progress-map-guided GAN generates the same subject's MRI at one- and ... & \cgap no* \\
SWoMo: Neuro-Symbolic World Model for Cataract Surgery Simulation & 3 & surgery/robotics & Explicit surgical tool-motion increments control tip position, orientation, articulation, ... & A rule-based Godot simulator advances geometry and tool-tissue interactions; a conditioned ... & \cgood yes \\
Treatment-aware Diffusion Probabilistic Model for Longitudinal MRI Generation ... & 3 & longitudinal progression & A target future treatment type and day are supplied; ... & Treatment-day embeddings condition a diffusion and segmentation network to generate the ... & \cgap no* \\
auton-survival: an Open-Source Package for Regression, Counterfactual Estimation, ... & 4 & treatment/causal & A treatment arm or exposure is supplied for counterfactual ... & The package fits treatment-specific survival functions and counterfactual survival regressors, with ... & \cgood yes \\
Double Robust Representation Learning for Counterfactual Prediction & 4 & treatment/causal & A binary treatment assignment indexes the two potential-outcome heads. & Double-robust representation learning combines treatment-specific outcome models with entropy-balancing weights that ... & \cgood yes \\
Agentifying Patient Dynamics within LLMs through Interacting with ... & 5 & treatment/causal & A settable 5-by-5 grid of intravenous-fluid and vasopressor-dose combinations. & A learned GRU clinical world model predicts the next physiological state ... & \cgood yes \\
ECG-WM: A Physiology-Informed ECG World Model for Clinical ... & 5 & treatment/causal & Drug identity, dose, and protocol drawn from a clinically ... & ODE-regularized latent diffusion simulates post-intervention ECG trajectories conditioned on drug action. & \cgap no* \\
LeJEPA + I-JEPA: A World-Model-Grounded Multi-Agent Framework for ... & outside & medical imaging & Prompt rewriting and report constraints are orchestration operations, not ... & No learned future-state transition is implemented; the released notebook labels LeJEPA ... & \cgood yes \\
A cost-optimized medical digital twin framework for secure ... & outside & other medical & ILP, PATI-Greedy, and HybridQeGA assign each computing task to ... & No learned or mechanistic patient, workflow, or infrastructure state transition is ... & \cgap no* \\
\bottomrule
\end{tabularx}
\vspace{2pt}
\\{\scriptsize Code: {\color{mwmteal}\faCheck} paper-linked public artifact; {\color{mwmamber}\faAdjust} partial/related artifact; {\color{mwmgray}\faMinus} no verified public study code or not applicable.}
\end{table*}

